%--------------------------------------------------------------------------
%		Não é necessário mudar qualquer coisa no preambulo deste arquivo
%		basta apenas colocar o seu texto, e as palavras-chave
%		observando as normas ABNT
%--------------------------------------------------------------------------
\chapter*{RESUMO} 
\thispagestyle{empty}

A irrigação responde por cerca de 70\% da água doce retirada no mundo, e o manejo feito por turnos fixos ou pela observação do produtor tende a aplicar mais ou menos água do que a cultura necessita. As soluções comerciais que automatizam essa decisão costumam ser proprietárias, ter custo elevado e oferecer pouca flexibilidade quanto aos sensores utilizados. Este trabalho teve como objetivo desenvolver um sistema de irrigação inteligente de baixo custo, baseado em Internet das Coisas, capaz de monitorar o solo e o ambiente, estimar a necessidade hídrica da cultura e automatizar a irrigação por gotejamento, com acompanhamento e controle a distância. Trata-se de uma pesquisa aplicada e exploratória, conduzida de forma incremental, em que um primeiro protótipo baseado em Raspberry Pi e MQTT deu lugar à arquitetura final, organizada em três camadas. Na camada de percepção, um dispositivo com ESP32, sensor de umidade do solo, sensor DHT11, módulo relé e display LCD lê os sensores e decide quando irrigar, com base em um limiar de umidade, e quanto irrigar, pelo método FAO-56, com a evapotranspiração de referência estimada pela equação de Hargreaves-Samani. Na camada de nuvem, uma API REST em Java com Spring Boot recebe os dados, aplica as regras de acesso, armazena as informações no Cloud Firestore e entrega os comandos ao dispositivo por long-polling. Na camada de aplicação, um painel web em Angular permite acompanhar as leituras e o histórico, configurar os parâmetros agronômicos e acionar a irrigação, com perfis de acesso distintos. O sistema foi avaliado em bancada por meio de dez cenários de teste de ponta a ponta, comprovados por registros do dispositivo, documentos do banco de dados, telas do painel e respostas da API. Os resultados mostraram que as leituras chegaram ao painel com os mesmos valores registrados no dispositivo, que os comandos foram executados nos modos manual e automático com os tempos calculados e que as travas de segurança e o controle de acesso funcionaram como previsto. Os testes também apontaram ajustes a fazer no firmware, como o registro do volume das irrigações interrompidas. Conclui-se que é viável construir, com componentes estimados entre R\$~200 e R\$~345, um sistema de irrigação automatizado e modular, embora a sua eficiência agronômica ainda dependa de uma avaliação em campo.

%--------------------------------------------------------------------------
\vspace{2cm}
{\Large \textbf{Palavras-chave:}}
\vspace{0.4cm}
%--------------------------------------------------------------------------

Internet das Coisas. Irrigação inteligente. ESP32. API REST. Evapotranspiração.
